\section{Relevância do trabalho}

A relevância do trabalho está associada à necessidade de modelar de forma integrada fenômenos orbitais e robóticos em operações de \gls{rvdb}.
A formulação desenvolvida permite analisar a influência direta do \gls{mr} na dinâmica da espaçonave perseguidora, aspecto que nem sempre é explicitamente considerado em modelos puramente orbitais, onde normalmente utilizam a \gls{hcw} para o movimento orbital relativo entre espaçonaves.

% A originalidade deste trabalho reside na integração, em um único ambiente de simulação, da dinâmica orbital relativa, da dinâmica acoplada espaçonave--\gls{mr} e das leis de controle aplicadas à operação do manipulador e à atitude da espaçonave, estruturadas sob uma formulação matematica que normalmente é utilizada de forma separada nos trabalhos associados.

Quanto à generalidade, as condições iniciais adotadas nos cenários simulados não restringem a aplicabilidade do modelo.
Elas foram selecionadas como hipóteses representativas para análise de desempenho, mas a estrutura matemática desenvolvida é válida para diferentes configurações orbitais e parâmetros do manipulador.

Do ponto de vista da aplicabilidade, o modelo pode servir como ferramenta preliminar para estudos de aproximação, inspeção, manutenção e atracação em órbita, contribuindo para análises conceituais e investigações acadêmicas sobre operações de \gls{rvdb}.